\documentclass[12pt]{article}

% --- Packages ---
\usepackage[utf8]{inputenc}
\usepackage[T1]{fontenc}
\usepackage[french]{babel}
\usepackage{amsmath, amssymb}
\usepackage{geometry}
\usepackage{tikz}
\usepackage[most]{tcolorbox}
\usepackage{multicol}
\usepackage{lmodern}
\usepackage{pgfplots}
\pgfplotsset{compat=1.18}
\usepackage{mdframed}
\usepackage{tkz-tab}
\usepackage{bm}
\usepackage{xcolor}
\usepackage{sectsty}
\usepackage[utf8]{inputenc}
\usepackage[T1]{fontenc}
\usepackage[inline]{enumitem}
\usepackage[dvipsnames]{xcolor}
\usepackage{multido}
% --- Page ---
\geometry{margin=2cm}

% Footer
\usepackage{fancyhdr}
\pagestyle{fancy}
\cfoot{\thepage}
\rfoot{\textit{G.M.}}
\renewcommand{\headrulewidth}{0pt}
\renewcommand{\footrulewidth}{0pt}

% Pour les normes : 
\newcommand*{\norme}[1]{\left\lVert\overrightarrow{#1}\right\rVert}

\usetikzlibrary{shapes.misc}

\tikzset{cross/.style={cross out, draw=black, minimum size=2*(#1-\pgflinewidth), inner sep=0pt, outer sep=0pt},
%default radius will be 1pt. 
cross/.default={4pt}}

% mettre dans le préambule, après \usepackage{xcolor}
\usepackage{sectsty} % gère facilement la couleur des titres
\setlength{\parindent}{0pt}

\sectionfont{\color{red}\Large\bfseries}     % pour \section*
\subsectionfont{\color{red}\large\bfseries\hspace{1cm}}

\begin{document}

% ============================
%   Fractions, puissances, racines carrées
% ============================

\newpage

\begin{center}
    {\LARGE \textbf{Fractions, puissances, racines carrées}}\\[4mm]
    {\large Seconde}
\end{center}

\hrule
\vspace{0pt}

% ============================
%   Partie 1 : Fractions
% ============================

\section*{\textcolor{red}{Partie 1 : Fractions}}

\subsection*{\hspace{1cm} 1) Calculs avec des fractions (rappels)}

\begin{mdframed}[
    linewidth=1pt,
    roundcorner=5pt,
    linecolor=Red,
    innertopmargin=10pt,
    innerbottommargin=10pt
]
{\color{red}
\underline{Propriété :}

\begin{center}
    $\dfrac{a}{D} + \dfrac{b}{D} = \dfrac{a+b}{D}$
    \hspace{1cm}$\dfrac{a}{D} - \dfrac{b}{D} = \dfrac{a-b}{D}$
    \hspace{1cm}$\dfrac{a}{b} \times \dfrac{c}{d} = \dfrac{a \times b}{c \times d}$
    \hspace{1cm}$\dfrac{a}{b} \div \dfrac{c}{d} = \dfrac{a}{b} \times \dfrac{d}{c}$
\end{center}
}
\end{mdframed}

\begin{mdframed}[
    linewidth=1pt,
    linecolor=RoyalBlue,
    innertopmargin=10pt,
    innerbottommargin=10pt
]
{\color{RoyalBlue}\textbf{\underline{Exercice d'application n°1 :} Effectuer des 
calculs de fractions}}

\begin{center}
    $A = \dfrac{5}{4}+\dfrac{6}{16}$
    \hspace{1cm}$B = \dfrac{5}{3}-\dfrac{6}{5}$
    \hspace{1cm}$C = \dfrac{2}{-3}\times\dfrac{-5}{11}$
    \hspace{1cm}$D = \genfrac{}{}{1pt}{0}{\phantom{\dfrac{1}{1}}\dfrac{3}{4}\phantom{\dfrac{1}{1}}}{\dfrac{-5}{8}}$
    \hspace{1cm}$E = \dfrac{8}{7}-\dfrac{4}{7} \times \dfrac{5}{3}$
\end{center}

\vspace{0.5cm}
\dotfill

\vspace{0.5cm}
\dotfill

\vspace{0.5cm}
\dotfill

\vspace{0.5cm}
\dotfill

\vspace{0.5cm}
\dotfill

\vspace{0.5cm}
\dotfill

\vspace{0.5cm}
\dotfill

\vspace{0.5cm}
\dotfill

\vspace{0.5cm}
\dotfill

\vspace{0.5cm}
\dotfill

\vspace{0.5cm}
\dotfill

\vspace{0.5cm}
\dotfill

\end{mdframed}

\vspace{10pt}

% ========================================================
%             2) Réduire des expressions au même dénominateur
% ========================================================
\subsection*{\hspace{1cm} 2) Réduire des expressions au même dénominateur}

\begin{mdframed}[
    linewidth=1pt,
    roundcorner=5pt,
    linecolor=Red,
    innertopmargin=10pt,
    innerbottommargin=10pt
]
{\color{red}
\underline{Propriété :}

\begin{center}
    $\dfrac{a}{{\color{blue}b}} + \dfrac{c}{{\color{ForestGreen}d}} = \dfrac{a{\color{ForestGreen}d}}{b{\color{ForestGreen}d}} + \dfrac{c{\color{blue}b}}{d{\color{blue}b}} = \dfrac{ad+cb}{bd}$
\end{center}
}
\end{mdframed}

\begin{mdframed}[
    linewidth=1pt,
    linecolor=RoyalBlue,
    innertopmargin=10pt,
    innerbottommargin=10pt
]
{\color{RoyalBlue}\textbf{\underline{Exercice d'application n°2 :} Réduire au même dénominateur}}

\begin{center}
    $A = \dfrac{7}{x-2} - \dfrac{5}{3}$ \hspace{2cm} $B = 3 + \dfrac{5x}{2x+1}$
\end{center}

\vspace{0.5cm}
\dotfill

\vspace{0.5cm}
\dotfill

\vspace{0.5cm}
\dotfill

\vspace{0.5cm}
\dotfill

\vspace{0.5cm}
\dotfill

\vspace{0.5cm}
\dotfill

\vspace{0.5cm}
\dotfill
\end{mdframed}


\vspace{10pt}
% ========================================================
%             Partie 2 : Puissances
% ========================================================
\section*{\textcolor{red}{Partie 2 : Puissances}}

% ========================================================
%             1) Rappels
% ========================================================
\subsection*{\hspace{1cm} 1) Rappels}

\begin{minipage}[t]{.45\textwidth}%
    \vspace{0pt}
    \begin{mdframed}[
        linewidth=1pt,
        roundcorner=5pt,
        linecolor=Red,
        innertopmargin=10pt,
        innerbottommargin=10pt
    ]
    {\color{red}
    \underline{Propriété :}

    \vspace{5pt}
    $a^4 = a \times a \times a \times a$

    \vspace{5pt}
    De façon générale : 

    \vspace{5pt}
    $a^n = \underbrace{a \times a \times a \times a \times a \times \text{\ ...\ } \times a}_{n \text{\ fois}}$
    }
    \end{mdframed}
\end{minipage}%
\hfill
\begin{minipage}[t]{.45\textwidth}%
    \vspace{0pt}
    \underline{Exemples :}

    \vspace{10pt}
    $2^3 =$

    \vspace{20pt}
    $11^5 = $

    \vspace{2cm}
\end{minipage}%

\begin{minipage}[t]{.45\textwidth}%
    \vspace{0pt}
    \begin{mdframed}[
        linewidth=1pt,
        roundcorner=5pt,
        linecolor=Red,
        innertopmargin=10pt,
        innerbottommargin=10pt
    ]
    {\color{red}
    \underline{Propriété :}

    \vspace{5pt}
    $a$ est un nombre non nul et $n$ un entier non nul.
    
    \vspace{5pt}
    \hspace{1cm}$a^1 = a$

    \hspace{1cm}$a^0 = 1$

    \hspace{1cm}$0^n = 0$

    \hspace{1cm}$1^n = 1$
    }
    \end{mdframed}
\end{minipage}%
\hfill
\begin{minipage}[t]{.45\textwidth}%
    \vspace{0pt}
    \underline{Exemples :}

    \vspace{20pt}
    $15^1 =$

    \vspace{20pt}
    $103^0 = $

    \vspace{20pt}
    $0^4 = $

    \vspace{20pt}
    $1^{12} = $
\end{minipage}%

\vspace{1cm}

{\color{red} \textbf{ATTENTION aux signes !}}

\vspace{10pt}
Ne pas confondre : 

\vspace{10pt}
$(-2)^4 = $ \hspace{6cm} et $-2^2 = $

\vspace{2cm}

% ========================================================
%             2) Opérations sur les puissances
% ========================================================
\subsection*{\hspace{1cm} 2) Opérations sur les puissances}

\vspace{5pt}
\begin{mdframed}[
    linewidth=1pt,
    roundcorner=5pt,
    linecolor=Red,
    innertopmargin=10pt,
    innerbottommargin=10pt
]
{\color{red}
\underline{Propriété :}

\vspace{5pt}
Avec $n$ et $p$ entiers relatifs :

\vspace{5pt}
\hspace{3cm} $a^n \times a^p = a^{n+p}$ \hspace{2cm} $\dfrac{a^n}{a^p} = a^{n-p}$ 
\hspace{2cm} $(a^n)^p = a^{n \times p}$

\vspace{10pt}
\hspace{3cm} $(a \times b)^n = a^n \times b^n$ \hspace{1.4cm} $a^{-1} = \dfrac{1}{a}$ 
\hspace{2.5cm} $a^{-n} = \dfrac{1}{a^n}$    
}
\end{mdframed}


\begin{mdframed}[
    linewidth=1pt,
    linecolor=RoyalBlue,
    innertopmargin=10pt,
    innerbottommargin=10pt
]
{\color{RoyalBlue}\textbf{\underline{Exercice d'application n°3 :} Effectuer des 
calculs sur les puissances}}

\begin{center}
    $A = \dfrac{1}{4^2}$
    \hspace{1cm}$B = 4^5 \times 4^7$
    \hspace{1cm}$C = \dfrac{5^4}{5^6}$
    \hspace{1cm}$D = 7^3 \times (7^2)^6$
    \hspace{1cm}$E = 6^7 \times 9^7$
\end{center}

\vspace{0.5cm}
\dotfill

\vspace{0.5cm}
\dotfill

\vspace{0.5cm}
\dotfill

\vspace{0.5cm}
\dotfill

\vspace{0.5cm}
\dotfill

\vspace{0.5cm}
\dotfill

\vspace{0.5cm}
\dotfill

\vspace{0.5cm}
\dotfill

\vspace{0.5cm}
\dotfill

\vspace{0.5cm}
\dotfill

\vspace{0.5cm}
\dotfill

\vspace{0.5cm}
\dotfill

\vspace{20pt}

{\color{RoyalBlue}\textbf{\underline{Exercice d'application n°3 :} Appliquer aux 
puissances de 10.}}

\vspace{5pt}
\hspace{1cm}a) Écrire sous la forme $10^n$ ou $10^{-n}$ :

\begin{center}
    $A = 10^4 \times 10^7$
    \hspace{1cm}$B = \dfrac{10^{-4}}{10^5}$
    \hspace{1cm}$C = (10^2)^{-6}$
    \hspace{1cm}$D = 10^{-4} \times (10^3)^{-1}$
\end{center}

\vspace{0.5cm}
{\color{black!30}\dotfill}

\vspace{0.5cm}
{\color{black!30}\dotfill}

\vspace{0.5cm}
{\color{black!30}\dotfill}

\vspace{0.5cm}
{\color{black!30}\dotfill}

\vspace{0.5cm}
{\color{black!30}\dotfill}

\vspace{0.5cm}
{\color{black!30}\dotfill}

\vspace{0.5cm}
{\color{black!30}\dotfill}

\vspace{0.5cm}
{\color{black!30}\dotfill}

\vspace{0.5cm}
{\color{black!30}\dotfill}

\vspace{0.5cm}
{\color{black!30}\dotfill}

\vspace{0.5cm}
{\color{black!30}\dotfill}

\newpage
\hspace{1cm}b) Écrire en notation scientifique :

\begin{center}
    $A = 4 \times 7 \times 10^{-5} \times 10^{-8}$
    \hspace{1cm}$B = \dfrac{7 \times 10^{-4} \times 5 \times 10^8}{56 \times 10^{-9}}$
    \hspace{1cm}$C = \dfrac{32 \times 10^{-4} + 6 \times 10^{-3}}{2 \times 10^{-5}}$
\end{center}

\vspace{0.5cm}
{\color{black!30}\dotfill}

\vspace{0.5cm}
{\color{black!30}\dotfill}

\vspace{0.5cm}
{\color{black!30}\dotfill}

\vspace{0.5cm}
{\color{black!30}\dotfill}

\vspace{0.5cm}
{\color{black!30}\dotfill}

\vspace{0.5cm}
{\color{black!30}\dotfill}

\vspace{0.5cm}
{\color{black!30}\dotfill}

\vspace{0.5cm}
{\color{black!30}\dotfill}

\vspace{0.5cm}
{\color{black!30}\dotfill}

\vspace{0.5cm}
{\color{black!30}\dotfill}

\vspace{0.5cm}
{\color{black!30}\dotfill}

\vspace{0.5cm}
{\color{black!30}\dotfill}

\vspace{0.5cm}
{\color{black!30}\dotfill}

\vspace{0.5cm}
{\color{black!30}\dotfill}

\vspace{20pt}


\end{mdframed}

\vspace{20pt}

% ========================================================
%             PARTIE 3 : Racines carrées
% ========================================================


\section*{\textcolor{red}{Partie 3 : Racines carrées}}

% ========================================================
%             1) Définition
% ========================================================
\subsection*{\hspace{1cm} 1) Définition}

\vspace{5pt}
\begin{mdframed}[
    linewidth=1pt,
    linecolor=ForestGreen,
    innertopmargin=10pt,
    innerbottommargin=10pt
]
{\color{ForestGreen}\textbf{\underline{Définition :}}}  

La racine carrée de $a$ est le nombre (\textbf{toujours positif}) dont le carré est $a$.

\end{mdframed}

\begin{mdframed}[
    linewidth=1pt,
    roundcorner=5pt,
    linecolor=Red,
    innertopmargin=10pt,
    innerbottommargin=0pt
]
{\color{red}
\textbf{\underline{Propriété :}}

\vspace{-15pt}
\begin{multicols}{3}
    $\sqrt{0} = 0$ \\
    $\sqrt{25} = 5$ \\
    $\sqrt{100} = 10$ 
\end{multicols}
\vspace{-25pt}
\begin{multicols}{3}
    $\sqrt{1} = 1$ \\
    $\sqrt{36} = 6$ \\
    $\sqrt{121} = 11$ 
\end{multicols}
\vspace{-25pt}
\begin{multicols}{3}
    $\sqrt{4} = 2$ \\
    $\sqrt{49} = 7$ \\
    $\sqrt{144} = 12$ 
\end{multicols}
\vspace{-25pt}
\begin{multicols}{3}
    $\sqrt{9} = 3$ \\
    $\sqrt{64} = 8$ \\
    $\sqrt{169} = 13$ 
\end{multicols}
\vspace{-25pt}
\begin{multicols}{3}
    $\sqrt{16} = 4$ \\
    $\sqrt{81} = 9$ \\
\end{multicols}

}
\end{mdframed}

\underline{Remarque :}

\vspace{5pt}
$\sqrt{-5} = $? \hspace{1cm} {\color{red}\fontencoding{U}\fontfamily{futs}\selectfont\char 49\relax} 
{\color{black!30}\dotfill}

\vspace{10pt}

\subsection*{\hspace{1cm} 2) Propriétés sur les racines}

\vspace{10pt}
\begin{mdframed}[
    linewidth=1pt,
    roundcorner=5pt,
    linecolor=Red,
    innertopmargin=10pt,
    innerbottommargin=0pt
]
{\color{red}
\textbf{\underline{Propriétés :}}

\vspace{5pt}
$a$ et $b$ deux nombres positifs.

\begin{center}
    $\sqrt{a} \times \sqrt{b} = \sqrt{a \times b}$ \hspace{1cm} 
    $\dfrac{\sqrt{a}}{\sqrt{b}} = \sqrt{\dfrac{a}{b}} \quad (b \neq 0)$ \hspace{1cm}
    $(\sqrt{a})^2 = a$ \hspace{1cm}
    $\sqrt{a^2} = a$
\end{center}

}
\end{mdframed}

{\color{red}\fontencoding{U}\fontfamily{futs}\selectfont\char 49\relax} De façon générale : 
$\sqrt{a} + \sqrt{b} \neq \sqrt{a + b}$ \hspace{1cm} et \hspace{1cm} $\sqrt{a} - \sqrt{b} \neq \sqrt{a-b}$


\vspace{10pt}

\begin{mdframed}[
    linewidth=1pt,
    roundcorner=5pt,
    linecolor=RoyalBlue,
    innertopmargin=10pt,
    innerbottommargin=10pt
]
{\color{RoyalBlue}
\textbf{\underline{Exercice d'application n°4:} Effectuer des calculs sur les racines carrées.}}

\vspace{5pt}
Écrire le plus simplement possible : 

$A = \sqrt{32} \times \sqrt{2}$ \hspace{1cm} $B = \sqrt{3} \times \sqrt{27}$ 
\hspace{1cm} $C = \sqrt{3} \times \sqrt{36} \times \sqrt{3}$

\vspace{5pt}
$D = \dfrac{\sqrt{32}}{\sqrt{2}}$ \hspace{1cm} $E = \dfrac{\sqrt{50}}{\sqrt{72}}$ 
\hspace{1cm} $F = (4\sqrt{5})^2$ \hspace{1cm} $G = \dfrac{\sqrt{32} \times \sqrt{10}}{\sqrt{80}}$

\vspace{0.5cm}
{\color{black!30}\dotfill}

\vspace{0.5cm}
{\color{black!30}\dotfill}

\vspace{0.5cm}
{\color{black!30}\dotfill}

\vspace{0.5cm}
{\color{black!30}\dotfill}

\vspace{0.5cm}
{\color{black!30}\dotfill}

\vspace{0.5cm}
{\color{black!30}\dotfill}

\vspace{0.5cm}
{\color{black!30}\dotfill}

\vspace{0.5cm}
{\color{black!30}\dotfill}

\vspace{0.5cm}
{\color{black!30}\dotfill}

\vspace{20pt}
\end{mdframed}

\vspace{10pt}

\subsection*{\hspace{1cm} 3) Extraire un carré parfait et simplifier une expression}

\vspace{5pt}
{\color{RedOrange}\textbf{\underline{Méthode :} Extraire un carré parfait.}}

\vspace{5pt}
Écrire sous la forme $a\sqrt{b}$, avec $a$ et $b$ entiers et $b$ étant le plus petit possible.

\vspace{5pt}
On commence par faire apparaître un carré parfait sous la racine :

\vspace{20pt}
\hspace{1cm}$\sqrt{32} = $

\vspace{2cm}

\begin{mdframed}[
    linewidth=1pt,
    roundcorner=5pt,
    linecolor=RoyalBlue,
    innertopmargin=10pt,
    innerbottommargin=10pt
]
{\color{RoyalBlue}
\textbf{\underline{Exercice d'application n°5:} Simplifier des racines.}}

\vspace{5pt}
Écrire sous la forme $a\sqrt{b}$, avec $a$ et $b$ entiers et $b$ étant le plus petit possible.

\vspace{5pt}
$A = \sqrt{72}$ \hspace{1cm} $B = \sqrt{45}$ \hspace{1cm} $C = 3\sqrt{125}$

\vspace{0.5cm}
{\color{black!30}\dotfill}

\vspace{0.5cm}
{\color{black!30}\dotfill}

\vspace{0.5cm}
{\color{black!30}\dotfill}

\vspace{0.5cm}
{\color{black!30}\dotfill}

\vspace{0.5cm}
{\color{black!30}\dotfill}

\vspace{0.5cm}
{\color{black!30}\dotfill}

\vspace{0.5cm}
{\color{black!30}\dotfill}

\vspace{0.5cm}
{\color{black!30}\dotfill}

\vspace{20pt}

\end{mdframed}

\vspace{10pt}

\subsection*{\hspace{1cm} 4) Simplifier les écritures complexes contenant des racines carrées}

\vspace{5pt}
{\color{RedOrange}\textbf{\underline{Méthode :} Simplifier une écriture}}

\vspace{5pt}
On regroue les membres de la même \textbf{« famille de racines irréductibles »} pour réduire l'expression.

\vspace{5pt}
Les différentes familles sont donc : $\sqrt{2}$, $\sqrt{3}$, $\sqrt{5}$, $\sqrt{6}$, $\sqrt{7}$... 

\vspace{20pt}
Par exemple : $4\sqrt{2} - 5\sqrt{2} = $

\vspace{2cm}

\begin{mdframed}[
    linewidth=1pt,
    roundcorner=5pt,
    linecolor=RoyalBlue,
    innertopmargin=10pt,
    innerbottommargin=10pt
]
{\color{RoyalBlue}
\textbf{\underline{Exercice d'application n°6:} Simplifier une écriture}}

\vspace{5pt}
\hspace{1cm}1) Écrire le plus simplement possible : 

\vspace{5pt}
$A = 4\sqrt{3} - 2\sqrt{3} + 6\sqrt{2}$ \hspace{1cm} $B = 7\sqrt{2} - 3\sqrt{5} + 8\sqrt{2} - \sqrt{5}$ 
\hspace{1cm} $C = (3 - 2\sqrt{3}) - (4 - 6\sqrt{3})$

\vspace{10pt}
\hspace{1cm}2) Écrire les expressions suivantes sous la forme $a\sqrt{b}$, où $a$ et $b$ sont des entiers 
et $b$ le plus petit possible.

\vspace{5pt}
$D = \sqrt{12} + 7\sqrt{3} - \sqrt{27}$ \hspace{1cm} $E = \sqrt{125} - 2\sqrt{20} + 6\sqrt{80}$

\vspace{0.5cm}
{\color{black!30}\dotfill}

\vspace{0.5cm}
{\color{black!30}\dotfill}

\vspace{0.5cm}
{\color{black!30}\dotfill}

\vspace{0.5cm}
{\color{black!30}\dotfill}

\vspace{0.5cm}
{\color{black!30}\dotfill}

\vspace{0.5cm}
{\color{black!30}\dotfill}

\vspace{0.5cm}
{\color{black!30}\dotfill}

\vspace{0.5cm}
{\color{black!30}\dotfill}

\vspace{0.5cm}
{\color{black!30}\dotfill}

\vspace{0.5cm}
{\color{black!30}\dotfill}

\vspace{0.5cm}
{\color{black!30}\dotfill}

\vspace{0.5cm}
{\color{black!30}\dotfill}

\vspace{0.5cm}
{\color{black!30}\dotfill}

\vspace{0.5cm}
{\color{black!30}\dotfill}

\vspace{0.5cm}
{\color{black!30}\dotfill}

\vspace{20pt}

\end{mdframed}

\vspace{10pt}
\subsection*{\hspace{1cm} 5) Racines carrées et développements}

\vspace{5pt}
{\color{RedOrange}\textbf{\underline{Méthode :} Effectuer des développements avec des racines}}

\vspace{5pt}
On cherche s'il y a une identité remarquable, si oui, on l'applique, sinon on développe normalement.

\vspace{20pt}
Exemple : $(\sqrt{2} - 3)^2 = $ 

\vspace{3cm}

Exemple : $(2 + \sqrt{3})(\sqrt{2} - 3) = $ 

\vspace{2cm}

\begin{mdframed}[
    linewidth=1pt,
    roundcorner=5pt,
    linecolor=RoyalBlue,
    innertopmargin=10pt,
    innerbottommargin=10pt
]
{\color{RoyalBlue}
\textbf{\underline{Exercice d'application n°7:} Effectuer des développements avec des racines.}}

\vspace{5pt}
Développer et réduire les expressions suivantes : 

\vspace{5pt}
\hspace{1cm}$A = (\sqrt{3} - 4)^2$ \hspace{3.15cm} $B = (3+\sqrt{5})^2$

\vspace{5pt}
\hspace{1cm}$C = (\sqrt{2} - \sqrt{5})(\sqrt{2}+ \sqrt{5})$ \hspace{1cm} $D = (3+\sqrt{3})(1-\sqrt{2})$


\vspace{0.5cm}
{\color{black!30}\dotfill}

\vspace{0.5cm}
{\color{black!30}\dotfill}

\vspace{0.5cm}
{\color{black!30}\dotfill}

\vspace{0.5cm}
{\color{black!30}\dotfill}

\vspace{0.5cm}
{\color{black!30}\dotfill}


\vspace{20pt}
\end{mdframed}

\end{document}
